% Datasets (draft, Week 12; written 4 Oct 2026 from docs/v2/data_sources.md).
% Counts come from numbers.tex (Week 23 notation pass).

\section{Data and failure model}
\label{sec:data}

We use three public datasets and no proprietary data. Each is used under its published terms;
Table~\ref{tab:data} summarises them.

\begin{table}[t]
  \centering
  \small
  \begin{tabular}{@{}llll@{}}
    \toprule
    Dataset & Period & Size & Used for \\
    \midrule
    Backblaze Drive Stats~\cite{backblaze2026drivestats} & Apr 2013--Jun 2026 & \VDriveDaysM{}M drive-days, \VDrives{} drives & failure model \\
    Google cluster trace~\cite{googleclusterdata2019,tirmazi2020borg} & May 2019 & 3 of 8 cells (+1), 10--12k machines each & empirical $\rho$ \\
    Azure Public Dataset V1~\cite{azurepublicdataset,cortez2017resource} & 2017 & 2{,}013{,}767 VMs & VM size, tier, lifetime \\
    \bottomrule
  \end{tabular}
  \caption{Datasets. Counts are after cleaning (Section~\ref{sec:data}).}
  \label{tab:data}
\end{table}

\paragraph{Drive failures.}
Backblaze publishes one record per drive per day, with the drive's model, capacity, a failure
flag and its SMART attributes~\cite{backblaze2026drivestats}. We converted all 45 published
archives (35.1\,GB compressed) to columnar files and checked every archive's row count against
its CSV files. Two problems in the source were found and handled: two daily files use a
different date format (2018-02-25 and 2019-06-17; the first also uses old Mac line endings), and the file named
2022-12-10 contains an exact copy of the next day, so that day is missing and the duplicate is
dropped. A drive fails on its first record with the failure flag; later records are discarded.
Drives that stop reporting without a failure flag are censored. Restricted to 2013--2017 and to
drives of at least 1.5\,TB, our rebuild reproduces the per-drive table used by our earlier
prototype: 122{,}620 drives against 122{,}619, and 6{,}907 failures against 6{,}913, with the
remaining differences traced to individual drives. Labels for the failure model mark a
drive-day positive if the drive fails within $H$ days, and exclude drive-days whose outcome is
censored (Section~\ref{sec:model}). A drive is a proxy for a host: a host's failure probability is
composed from its drives and a separately assumed non-disk term (Section~\ref{sec:sim}).

\paragraph{Placement.}
Google's 2019 trace records eight Borg cells for 31 days, including every task's priority and
scheduling class and every machine's add, remove and update events~\cite{tirmazi2020borg}. We
use three cells chosen before any data was read, the ones Tirmazi et al.\ report as having the
most production, best-effort batch and mid-tier work, to span the range of workload mixes. We
map priority tiers to the three criticality tiers of the Azure data (production and monitoring
to high, mid-tier to medium, free and best-effort batch to low) and test two alternative
mappings. The trace has no hardware-failure labels. Our risk proxy is a machine platform's rate
of \emph{long} removals (a removal not followed by re-addition within 24 hours), because short
removals are dominated by routine maintenance~\cite{tirmazi2020borg}; this is a proxy, not a
failure rate, and every figure that uses it says so.

\paragraph{Virtual machines.}
The Azure Public Dataset V1 gives, for about two million VMs observed over 30 days in 2017,
their category (interactive, delay-insensitive or unknown), size and CPU
utilisation~\cite{cortez2017resource}. We use the category as the criticality tier and sample
VM size, CPU, tier and lifetime jointly from it. VMs in this trace are short-lived (median
lifetime 40 minutes), which matters for what a migration can protect
(Section~\ref{sec:problem}).

\paragraph{Failure model.}
\label{sec:model}
We predict whether a drive fails within $H$ days from its SMART attributes, their recent trends,
its model and its age, and train only on the past: training ends in August 2025, the validation
period (used to choose the model and fit its calibration) ends in January 2026, and the test
period runs to June 2026. At $H=30$ days a gradient-boosted model reaches a test PR-AUC of
\VPrAucSMART{} and ROC-AUC of \VRocAucSMART; the prototype's features (model, capacity and age),
refitted on the same split, reach a PR-AUC of only \VPrAucVOneFeatures, barely above the base
rate. The signal is in the SMART telemetry, not the model class: a logistic regression on the
same features does as well (\VPrAucLogit). We chose the predictor for the simulator by a rule
fixed in advance, the highest validation PR-AUC, which selected an ensemble of boosted models
(\VPrAucEnsembleVal{} against \VPrAucLogitVal{} for the next best). Isotonic calibration fitted on
the validation period makes its probabilities usable as $p$ in Eq.~\eqref{eq:benefit}. Even so,
most drives the model flags do not fail within the horizon, which is why what a migration costs
matters.

\paragraph{Where real placement sits.}
\label{sec:rho}
For each cell we compute Kendall's $\tau_b$ between the criticality tier of every placement and
the risk proxy of its machine, over three weekly windows after a one-week warm-up, with
bias-corrected bootstrap intervals over machines. To compare with the simulator, each value is
also mapped to the latent axis of a Gaussian copula under that cell's own margins ($\rho_L$). The
method, the three cells and the decision rule were fixed before any data was read. In all three
cells, under every combination of three tier mappings, three proxies and two weightings,
$|\tau_b| \le \VGOneMaxAbsTauB$ and the upper end of every latent interval is at most
\VGOneMaxLatentUpper. The sign of the estimate changes with the analysis choices; its magnitude
does not. A fourth cell, chosen by a rule fixed before downloading it and analysed separately,
gives $\tau_b = \VSecondTauB$ ($\rho_L = \VSecondLatent$). Critical work is essentially uncorrelated
with machine risk, so the realistic regimes of our evaluation (B1 and its interval ends) place VMs
at these measured values.
